\documentclass{reportkit}

\setreportkitleftheader{REPORTKIT / ALGORITHM VISUALIZATION TEST}
\setreportkitfooter{Algorithm visualization -- linear containers acceptance test}
\setreportkitversion{v1.10.0}

\title{ReportKit -- linear-container algorithm-state primitives}
\author{Test build}
\date{17 September 2026}

\begin{document}
\maketitle

\begin{execsummary}
This is an acceptance test for reportkit-algorithm-linear.sty's linear
container primitives: \texttt{stackstate} (LIFO) and \texttt{queuestate}
(FIFO), the first two members of primitive family C
(docs/superpowers/specs/2026-09-16-reportkit-algorithm-visualization-primitives-spec.md,
section 7). Both share one internal linear-container renderer and reuse the
existing \texttt{arraystate} cell and pointer chrome; neither introduces new
theme tokens.
\end{execsummary}

\section{Stack: bracket matching}
A \texttt{stackstate} tracks unmatched opening brackets while scanning a
string left to right. The most recently pushed bracket renders at the
visual top, explicitly marked with a TOP arrow, so the reader can see which
bracket the next closing character must match without relying on color
alone.

\begin{diagram}[type=stack,width=.5\textwidth,caption={Unmatched opening brackets while scanning \texttt{"([<)"}.},description={A stack holding three unmatched opening brackets -- round, square, and angle -- with the most recently pushed angle bracket marked current and the stack top labelled.}]
\begin{stackstate}
  \push{(}
  \push{[}
  \push[state=current]{<}
\end{stackstate}
\end{diagram}

\section{Stack: full state vocabulary}
All nine shared algorithm states applied to stack entries, so state meaning
stays distinguishable through border weight and line style, not fill color
alone.

\begin{diagram}[type=stack,width=.6\textwidth,caption={Every shared algorithm state applied to one stack entry each.},description={Nine stacked cells, one per shared algorithm state: unseen at the bottom through current at the top.}]
\begin{stackstate}
  \push[state=unseen]{uns}
  \push[state=blocked]{blk}
  \push[state=discarded]{disc}
  \push[state=resolved]{res}
  \push[state=visited]{vis}
  \push[state=frontier]{front}
  \push[state=candidate]{cand}
  \push[state=active]{act}
  \push[state=current]{cur}
\end{stackstate}
\end{diagram}

\section{Queue: BFS work queue}
A \texttt{queuestate} models the frontier of a breadth-first search. Items
render left to right in enqueue order, with the front (dequeue end, next
node to visit) marked on the left and the rear (enqueue end, where newly
discovered neighbours join) marked on the right.

\begin{diagram}[type=queue,width=\textwidth,caption={BFS frontier after visiting the start node's neighbours.},description={A three-element FIFO work queue with the front node marked current, the dequeue end labelled on the left, and the enqueue end labelled on the right.}]
\begin{queuestate}
  \enqueue[state=current]{B}
  \enqueue[state=frontier]{C}
  \enqueue[state=frontier]{D}
\end{queuestate}
\end{diagram}

\section{Queue: Kahn's algorithm}
The same primitive models a topological-sort work queue: nodes whose
in-degree has just reached zero are enqueued (frontier), while the node
currently being removed and processed sits at the dequeue end.

\begin{diagram}[type=queue,width=\textwidth,caption={Nodes with zero remaining in-degree, ready to process.},description={A four-element FIFO work queue for Kahn's algorithm, with the next node to process marked current on the left and newly ready nodes marked frontier on the right.}]
\begin{queuestate}
  \enqueue[state=current]{A}
  \enqueue[state=frontier]{E}
  \enqueue[state=frontier]{F}
  \enqueue[state=frontier]{G}
\end{queuestate}
\end{diagram}

\section{Grayscale and semantic check}
\begin{principle}{State is not just color}
Both primitives reuse \texttt{arraystate}'s shared \texttt{rk algo cell}
styling, so every state carries the same distinct line-style fragment (bold
solid for \texttt{current}, dashed for \texttt{candidate}, dotted for
\texttt{blocked}, and so on) as the array and window primitives. The TOP,
DEQUEUE, and ENQUEUE markers reuse the existing arraystate pointer chrome
rather than introducing a new visual language.
\end{principle}

\end{document}
